\scanpage{23}
\uline{Corollaire 2} Soit $A$ un élément normal de $\underline b$, $f$ une fonction borélienne sur $\underline C$ bornée sur tout compact, nulle au voisinage de $0$. Alors $f(A)\in\underline a$.

Soit maintenant $A$ un élément normal de $\underline b$. $A$ est contenu dans une sous-algèbre \inserted{auto-adjointe} commutative \unsure{uniformt} fermée $\underline c$ de $\underline b$, qui s'identifie à un $C_0(M)$, sur $M$ il y a une mesure $\unsure{\mu|\rho}$. \deleted{$\unsure{\mu_A}$} La mesure spectrale $\unsure{\mu_f|\nu_f}$ de $f$ existe, elle ne dépend pas du choix de $\underline c$, car on a pour toute $\varphi\in\mathcal K(C^*)$ :
\[
\langle\varphi,\unsure{\mu_f|\nu_f}\rangle=\int_M\varphi(f)=T(\varphi(A)).
\]
On l'appelle \uline{mesure spectrale de $A$}, et on la note $\unsure{\mu_A|\nu_A}$ :
\[
(1)\qquad\langle\varphi,\unsure{\mu_A|\nu_A}\rangle=T(\varphi(A))\qquad\varphi\in\mathcal K(C^*)
\]
\deleted{Si $\unsure{\varphi|T}$ est fidèle, le}
\inserted{La mesure spectrale de $A$ s'identifie à la mesure sur le spectre de l'algèbre auto-adjointe commutative \unsure{uniformt} fermée engendrée par $A$, \deleted{\unsure{identifiée}} définie \deleted{\unsure{dans} $(A)$} par le corollaire 1. En particulier, on a}
\[
(2)\qquad\boxed{\mathop{\rm supp}\unsure{\mu_A|\nu_A}=\unsure{\sigma^*(A)|\sigma(A)}}\qquad\text{(spectre de $A$).}
\]
Cette inclusion est une égalité, quand $\unsure{T|\varphi}$ est fidèle. \unsure{La méthode montre} que $A\in\underline a\Longleftrightarrow\unsure{z}$ sommable pour $\unsure{\mu_A|\nu_A}$, \deleted{Alors on a \unsure{$\langle z^n,\mu_A\rangle$}}
\[
(3)\qquad\int z^n\,d\unsure{\mu_A|\nu_A}(z)=T(A^n)\qquad(n\geq1)
\]

Supposons maintenant $A\geq0$. Alors \deleted{$\unsure{A}$} $f$ admet un réarrangement décroissant, qui ne dépend que de $A$ (car $\unsure{\mu_A|\nu_A}$ ne dépend que de $A$), noté $\varphi_A$, et appelé \uline{réarrangement spectral de $A\in\underline b^+$}. C'est une fonction \inserted{bornée} décroissante sur $R^+_*$, $\|\varphi_A\|=\|A\|$. On pose de plus $\Phi_A=\Phi_f$, $\Psi_A=\Psi_f$, (pour $A$ positif)
\deleted{$\Phi_A=\Phi_f$, $\Psi_A=\Psi_f$}
\[
(4)\qquad\deleted{\left\{\begin{aligned}\Phi_A&=\Phi_f\\\Psi_A&=\Psi_f\\\Delta_A&=\Delta_f=\exp\Psi_f\end{aligned}\right.}
\quad
\left\{\begin{aligned}\Phi_A(t)&=\int_0^t\varphi_A(s)\,ds\\\Psi_A(t)&=\int_0^t\log\varphi_A(s)\,ds\end{aligned}\right.
\]
\deleted{$\Delta_A(t)=\exp\int_0^t\log\varphi_A(s)\,ds$}
\mspage{24}
